% !TeX spellcheck = ru_RU
% !TEX root = vkr.tex

\section{Особенности реализации}

В данном разделе будут описаны ключевые особенности реализации статического анализатора.

\subsection{Архитектура}

Статический анализатор разделён на два модуля.
Первый модуль отвечает за анализ декларативной спецификации из LLVM, второй ответственен за анализ модели на Sail.

\begin{itemize}
  \item Модуль, ответственный за анализ LLVM, получает на вход скомпилированные в единый \texttt{JSON} \texttt{.td} файлы, содержащие описание спецификации RISC-V ISA.
        Затем из входного \texttt{JSON} файла для каждой инструкции извлекаются следующие свойства:
        \begin{itemize}
          \item мнемоника (извлекается по ключу \texttt{AsmString});
          \item входные операнды (по ключу \texttt{InOperandList});
          \item выходные операнды (по ключу \texttt{OutOperandList});
          \item неявные входные операнды (по ключу \texttt{Uses});
          \item неявные выходные операнды (по ключу \texttt{Defs});
          \item читает ли инструкция память (по ключу \texttt{mayLoad});
          \item пишет ли инструкция в память (по ключу \texttt{mayStore});
          \item для какой разрядности (\texttt{RV32} или \texttt{RV64}, или для обеих) определена инструкция (по ключу \texttt{Predicates}).
        \end{itemize}
        Наконец, анализатор генерирует модуль, который содержит всю собранную информацию о свойствах каждой инструкции и предоставляет интерфейс для доступа к информации.
        Генерация отдельного модуля со всей информацией реализована для того, чтобы при каждом запуске анализатора не анализировать один и тот же \texttt{JSON}, тем самым экономя драгоценное время.

  \item Второй модуль, отвечающий за статический анализ модели на Sail, получает на вход \texttt{.sail} файлы, которые содержат определение модели.
        Далее при помощи библиотеки компилятора языка Sail анализатор производит синтаксический анализ и вывод типов, на выходе получая типизированное AST.
        Затем на полученном ранее типизированном AST запускаются анализы, о которых подробнее далее.
\end{itemize}

\subsection{Специализация модели}

Как было сказано ранее эмулятор sail-riscv представляет из себя интерпретатор некоторого AST, каждый узел которого параметризован некоторыми параметрами.
Один узел AST модели может представлять сразу несколько инструкций, но специализация узла представляет из себя конкретную инструкцию.

Например, узел AST из листинга \ref{itype} специализированный по последнему параметру
\begin{lstlisting}[caption={}, language={}]
  ITYPE (imm, rs1, rd, RISCV_ADDI)
\end{lstlisting}
представляет инструкцию
\begin{lstlisting}[caption={}, language={}]
  addi rd, rs1, imm
\end{lstlisting}

\subsection{Сопоставление узлов AST и мнемоник}

Во-первых, необходимо было определить за какие инструкции отвечает каждая специализация узла AST в модели.
Для этого производился анализ специализаций функций \texttt{assembly}, которые как раз отвечают за отображение узлов AST в мнемоники инструкций.

\subsection{Поиск входных и выходных операндов}

В модели sail-riscv для взаимодействия с регистрами используются следующие функции:
\begin{itemize}
  \item \texttt{rX} для чтения, \texttt{wX} для записи GPR регистров;
  \item \texttt{rF} для чтения, \texttt{wF} для записи вещественных регистров;
  \item \texttt{rV} для чтения, \texttt{wV} для записи векторных регистров.
\end{itemize}

Для поиска входных и выходных операндов строился граф вызовов всех функций для всех специализаций функций \texttt{execute}, которые отвечают за интерпретацию поведения инструкций.
Затем производился обход этого графа вызовов и определение того, какие функции с какими аргументами вызывали функции для взаимодействия с регистрами.

Дополнительно для каждой функции производился анализ алиасов\footnote{\url{https://en.wikipedia.org/wiki/Alias_analysis} (дата обращения: \DTMdate{2024-12-17}).}.
\begin{lstlisting}[caption={Пример функции, которая интерпретирует выполнение инструкции c.addi4spn}, language=Caml, frame=single, label=alias]
function clause execute (C_ADDI4SPN(rdc, nzimm)) = {
      let imm : bits(12) = (0b00 @ nzimm @ 0b00);
      let rd = creg2reg_idx(rdc);
      execute(ITYPE(imm, sp, rd, RISCV_ADDI))
}
\end{lstlisting}
Если не производить анализ алиасов, то, например, на листинге \ref{alias} входным регистром был бы регистр \texttt{rd}, но на самом деле входным регистром является регистр \texttt{rdc}.

\subsection{Анализ обращения инструкции к памяти}

Для работы с памятью в sail-riscv используются функции \texttt{read\_ram} и \texttt{write\_ram}.
Для того чтобы определить читает или пишет инструкция память, для каждой специализации функции \texttt{execute} проверялась достижимость вершин \texttt{read\_ram} и \texttt{write\_ram} в ранее определённом графе вызовов всех функций.

\textbf{тут дичь. поправить}

Однако любая инструкция, которая работает с адресами, для трансляции виртуального адреса в физический обращалась к таблице страниц.
И при чтении таблицы страниц могла происходить как запись в память, так и чтение из памяти.
\begin{itemize}
  \item Чтение из памяти могло происходить в случае, если при трансляции виртуального адреса в физический происходил TLB miss. Тогда необходимо было прочитать таблицу страниц из памяти.
  \item Запись в память могла происходить тогда, когда при трансляции виртуального адреса в физический необходимо было обновить бит A или D у записи таблицы страниц.
\end{itemize}

По требованию заказчика обращения к памяти при трансляции адресов игнорировались. Для этого при решении задачи достижимости вершин \texttt{read\_ream} и \texttt{write\_ram} в графе вызовов игнорировался вызов функции \texttt{translateAddr}, отвечающей за трансляцию виртуального адреса в физический.

\subsection{Особенности расширения <<P>> в sail-riscv-p}

Проект sail-riscv-p\footnote{\url{https://github.com/umcann123/sail-riscv-p} (дата обращения: \DTMdate{2024-12-17}).} --- это форк эмулятора sail-riscv с реализованным расширением <<P>> (packed SIMD).

\begin{lstlisting}[caption={Пример функции \texttt{execute}, которая может интерпретировать регистры как парные}, language=Caml, frame=single, label=depxlen]
function clause execute (PEXT_64BITS_ADD_SUB(rs2, rs1, rd, op))
= {
  results : bits(64) = zeros();
  rs1_val : bits(64) = zeros();
  rs2_val : bits(64) = zeros();
  if sizeof(xlen) == 32
  then {
    rs1_val = X(rs1 | 0b00001) @ X(rs1 & 0b11110);
    rs2_val = X(rs2 | 0b00001) @ X(rs2 & 0b11110);
  }
  else {
    rs1_val = X(rs1);
    rs2_val = X(rs2);
  };
  results = match op {
    ADD64   => rs1_val + rs2_val,
    ... \\ etc
  };
  if sizeof(xlen) == 32
  then {
    X(rd | 0b00001) = results[63..32];
    X(rd & 0b11110) = results[31..0];
  }
  else X(rd) = results;
  RETIRE_SUCCESS
}
\end{lstlisting}

На листинге \ref{depxlen} видно, что функция \texttt{PEXT\_64BITS\_ADD\_SUB} зависит от \texttt{xlen} --- разрядности RISC-V, и в зависимости от \texttt{xlen} функция \texttt{execute} работает с регистрами по разному.
В случае если \texttt{xlen} равен 32, то входные и выходные регистры интерпретируются как парные, иначе, если \texttt{xlen} не равно 32, то входные и выходные регистры интерпретируются как обычно.

Чтобы определить интерпретирует ли инструкиция регистры как парные, производилось приведение программы к ANF\footnote{\url{https://en.wikipedia.org/wiki/A-normal_form} (дата обращения: \DTMdate{2024-12-17})} форме, а затем строился и анализировался data flow graph.

Во-первых, для того чтобы аргументами всех функций были только тривиальные значения (константы и переменные) было реализовано приведение представления программы к ANF форме.
\begin{lstlisting}[caption={Вызов чтения из регистра до ANF}, language={Caml}, frame=single]
X(rs1 | 0b00001)
\end{lstlisting}

\begin{lstlisting}[caption={Вызов чтения из регистра после ANF}, language={Caml}, frame=single]
let a1 = rs1 | 0b00001 in X(a1)
\end{lstlisting}

Затем анализировался data flow graph каждой нормализованной \texttt{execute} функции.


\subsection{CI}

Был добавлен CI, который проверяет успешность сборки анализатора и проход тестов.

\section{Найденные расхождения}

Для LLVM версии 19 и sail-riscv\footnote{\url{https://github.com/riscv/sail-riscv/commits/05b845c91d1c1db7b361fc8d06e815b54ca0b07a} (дата обращения: \DTMdate{2025-05-17}).} при помощи анализатора были найдены следующие расхождения:

\begin{enumerate}
  \item \texttt{sp} не является явным аргументом в sail-riscv для инструкций:
        \texttt{c.addi16sp}, \texttt{c.addi4spn}, \texttt{c.fldsp}, \texttt{c.flwsp}, \texttt{c.fsdsp}, \texttt{c.fswsp}, \texttt{c.ldsp}, \texttt{c.lwsp}, \texttt{c.sdsp}, \texttt{c.swsp}.

        % \item \texttt{Zreg} является неявным входным регистром в Sail для инструкций:
        % \texttt{c.beqz}, \texttt{c.bnez}, \texttt{c.fldsp}, \texttt{c.li}, \texttt{c.mv}.

        % \item \texttt{Zreg} является неявным выходным регистром в Sail для инструкций:
        % \texttt{c.j}, \texttt{c.jr}.

  \item \texttt{RA} является неявным out регистром в sail-riscv для инструкций:
        \texttt{c.jal}, \texttt{c.jalr}

  \item В sail-riscv \texttt{vd} --- входной регистр для инструкций:
        \texttt{vadc.vim} \texttt{vadc.vvm} \texttt{vadc.vxm} \texttt{vcompress.vm} \texttt{vfmerge.vfm} \texttt{vmadc.vi} \texttt{vmadc.vim} \texttt{vmadc.vv} \texttt{vmadc.vvm} \texttt{vmadc.vx} \texttt{vmadc.vxm} \texttt{vmand.mm} \texttt{vmandn.mm} \texttt{vmerge.vim} \texttt{vmerge.vvm} \texttt{vmerge.vxm} \texttt{vmnand.mm} \texttt{vmnor.mm} \texttt{vmor.mm} \texttt{vmorn.mm} \texttt{vmsbc.vv} \texttt{vmsbc.vvm} \texttt{vmsbc.vx} \texttt{vmsbc.vxm} \texttt{vmxnor.mm} \texttt{vmxor.mm} \texttt{vsbc.vvm} \texttt{vsbc.vxm}.

  \item В sail-riscv \texttt{vd} и (неявно) \texttt{v0} --- входные регистры для инструкций:
        \texttt{vaaddu.vv} \texttt{vaaddu.vx} \texttt{vaadd.vv} \texttt{vaadd.vx} \texttt{vadd.vi} \texttt{vadd.vv} \texttt{vadd.vx} \texttt{vand.vi} \texttt{vand.vv} \texttt{vand.vx} \texttt{vasubu.vv} \texttt{vasubu.vx} \texttt{vasub.vv} \texttt{vasub.vx} \texttt{vdivu.vv} \texttt{vdivu.vx} \texttt{vdiv.vv} \texttt{vdiv.vx} \texttt{vfadd.vf} \texttt{vfadd.vv} \texttt{vfclass.v} \texttt{vfcvt.f.xu.v} \texttt{vfcvt.f.x.v} \texttt{vfcvt.rtz.x.f.v} \texttt{vfcvt.rtz.xu.f.v} \texttt{vfcvt.x.f.v} \texttt{vfcvt.xu.f.v} \texttt{vfdiv.vf} \texttt{vfdiv.vv} \texttt{vfirst.m} \texttt{vfmacc.vf} \texttt{vfmacc.vv} \texttt{vfmadd.vf} \texttt{vfmadd.vv} \texttt{vfmax.vf} \texttt{vfmax.vv} \texttt{vfmin.vf} \texttt{vfmin.vv} \texttt{vfmsac.vf} \texttt{vfmsac.vv} \texttt{vfmsub.vf} \texttt{vfmsub.vv} \texttt{vfmul.vf} \texttt{vfmul.vv} \texttt{vfmv.s.f} \texttt{vfmv.v.f} \texttt{vfncvt.f.f.w} \texttt{vfncvt.f.xu.w} \texttt{vfncvt.f.x.w} \texttt{vfncvt.rod.f.f.w} \texttt{vfncvt.rtz.x.f.w}...\footnote{Полный список разхождений --- \url{https://github.com/s-khechnev/llvm_td_lint_playground/blob/ins_outs_regs/Results.md} (дата обращения: \DTMdate{2025-03-12}).}.


\end{enumerate}

Для реализации расширения <<P>> в LLVM и в sail-riscv-p\footnote{\url{https://github.com/umcann123/sail-riscv-p} (дата обращения: \DTMdate{2024-12-17}).} при помощи анализатора были найдены следующие расхождения:

\begin{itemize}
  \item В sail-riscv-p \texttt{rd} является входным регистром для инструкций:
        \texttt{ave}, \texttt{bitrev}, \texttt{kaddw}, \texttt{kdmbb}, \texttt{kdmbb16}, \texttt{kdmbt}, \texttt{kdmbt16}, \texttt{kdmtt}, \texttt{kdmtt16}, \texttt{khmbb16}, \texttt{khmbt16}, \texttt{khmtt16}, \texttt{kmda}, \texttt{kmmwb2}, \texttt{kmmwb2.u}, \texttt{kmmwt2}, \texttt{kmmwt2.u}, \texttt{kmxda}, \texttt{ksllw}, \texttt{kslraw}, \texttt{kslraw.u}, \texttt{ksubw}, \texttt{kwmmul}, \texttt{kwmmul.u}, \texttt{smbb16}, \texttt{smbt16}, \texttt{smdrs}, \texttt{smds}, \texttt{smmul}, \texttt{smmul.u}, \texttt{smmwb}, \texttt{smmwb.u}, \texttt{smmwt}, \texttt{smmwt.u}, \texttt{smtt16}, \texttt{smxds}, \texttt{sra.u}, \texttt{ukaddw}, \texttt{uksubw}
\end{itemize}



% Реализация в широком смысле: что таки было сделано.
% Часто это на самом деле несколько отдельных разделов, по одному на каждую \enquote{реализационную} задачу из постановки.
% Например, \enquote{Архитектура} и \enquote{Особенности реализации}, либо по разделу на каждую крупную подзадачу.
% Если разделы получились недостаточно эпичными (меньше пары страниц), сделайте их подразделами одного большого раздела, как тут.

% Если программной реализации нет, тут пишутся Ваши теоретические наработки, если есть, хоть в каком-то виде, то опишите, даже если серьёзно кодом надо будет заняться только в следующей части.

% Если работа на текущем этапе предполагает \emph{только} обзор, то
% \begin{itemize}
%     \item но всё равно же надо было попробовать воспроизвести чужие результаты, напишите про это сюда;
%     \item если нет, то, наверное, обзор получился годным и ценным сам по себе, источников 40--50 было рассмотрено%
%           \footnote{Это не шутка, в хороших работах, где целый семестр делался только обзор, оно примерно так и выходит.},
%           тогда ладно, раздел \enquote{Описание решения} можно не делать.
% \end{itemize}

% Для понимания того, как отчёт по учебной практике должен писаться, можно посмотреть видео ниже.
% Они про научные доклады и написание научных статей.
% Учебные практики и ВКР отличаются тем, что тут есть требования отдельных глав (слайдов) со списком задач и списком результатов.
% Но даже если Вы забьёте на требования, специфичные для ВКР, и соблюдете все рекомендации ниже, получившиеся скорее всего будет лучше, чем первая попытка 99\% ваших однокурсников.

% \begin{enumerate}
%     \item \enquote{Как сделать великолепный научный доклад} от Саймона Пейтона Джонса~\cite{SPJGreatTalk} (на английском).
%     \item \enquote{Как написать великолепную научную статью} от Саймона Пейтона Джонса~\cite{SPJGreatPaper} (на английском).
%     \item \enquote{Как писать статьи так, чтобы люди их смогли прочитать} от Дэрэка Драйера~\cite{DreyerYoutube2020} (на английском).
%           По предыдущим ссылкам разбирается, что должно быть в статье, т.е. как она должна выгля\-деть на высоком уровне.
%           Тут более низкоуровнево изложено, как должны быть устроены параграфы и т.п.
%     \item Или на русском от С.~Григорьева~\cite{SemenI,SemenII}.
%     \item Ещё можно посмотреть How to Design Talks~\cite{JhalaYoutube2020} от Ranjit Jhala, но мы думаем, что первых ссылок всем хватит.
% \end{enumerate}

% Можно глянуть примеры хороших работ прошлых лет на сайте кафедры системного программирования%
% \footnote{Сайт кафедры СП, архив практик и ВКР: URL: \url{https://se.math.spbu.ru/theses.html} (дата обращения: \DTMdate{2024-01-15}).}.
% Например, работа Москаленко Ивана Николаевича за 3-й курс%
% \footnote{\enquote{Создание базы данных изображений для глобальной локализации в пространстве}, URL: \url{https://se.math.spbu.ru/thesis_download?thesis_id=1199} (дата обращения: \DTMdate{2024-01-15}).}.

% \subsection{Первая задача}
% \label{subsec:task1}

% \subsection{Вторая задача}
% \label{subsec:task2}

% \subsection{Третья задача}
% \label{subsec:task3}

% \subsection{Некоторые типичные ошибки}
% Здесь мы будем собирать основные ошибки, которые случаются при написании текстов.
% В интернетах тоже можно найти коллекции типич\-ных косяков%
% \footnote{\href{https://www.read.seas.harvard.edu/~kohler/latex.html}{https://www.read.seas.harvard.edu/\textasciitilde kohler/latex.html} (дата обращения: \DTMdate{2022-12-16}).}.

% Рекомендуется по-умол\-ча\-нию использовать красивые греческие бук\-вы $\sigma$  и $\phi$, а именно $\phi$ вместо $\varphi$.
% В данном шаблоне команды для этих букв переставлены местами по сравнению с ванильным \TeX'ом.

% Также, если работа пишется на русском языке, необходимо, чтобы работа была написана на \textit{грамотном} русском языке даже если автор из ближнего зарубежья%
% \footnote{Теоретически, возможен вариант написания текстов на английском языке, но это необходимо обсудить в первую очередь с научным руководителем.}.
% Это включает в себя:
% \begin{itemize}
%     \item разделы должны оформляться с помощью \verb=\section{...}=, а также \verb=\subsection= и т.~п.; не нужно пытаться нумеровать вручную;
%     \item точки после окончания предложений должны присутствовать;
%     \item пробелы после запятых и точек, в конце слова перед скобкой;
%     \item неразрывные пробелы там, где нужны пробелы, но переносить на другую строку нельзя, например \verb=т.~е.=, \verb=А.~Н.~Терехов=, \verb=что-то~\cite{?}=, \verb=что-то~\ref{?}=;
%     \item дефис там, где нужен дефис (обозначается с помощью одиночного \enquote{минуса} (англ. dash) на клавиатуре);
%     \item двойной дефис там, где он нужен; а именно  при указании проме\-жутка в цифрах: в 1900--1910 г.~г., стр. 150--154;
%     \item тире (т.~е. \verb=---=~--- тройной минус) на месте тире, а не что-то другое; в русском языке тире не может \enquote{съезжать} на новую строку, поэтому стоит использовать такой синтаксис: \verb=До~--- после=;
%     \item даты стоит писать везде одинаково; чтобы об этом не следить, можно пользоваться заклинанием \verb=\DTMdate{2022-12-16}=;
%     \item правильные кавычки должны набираться с помощью пакета \texttt{csquotes}: для основного языка в \texttt{polyglossia} стоит использовать команду \verb=\enquote{текст}=, для второго языка стоит использовать \verb=\foreignquote{язык}{текст}=; правильные кавычки в русской типографии~--- \verb=<<ёлочки>>=, ни в коем случае не \verb="скандинавские лапки"=;
%     \item все перечисления должны оформляться с помощью \verb=\enumerate= или \verb=\itemize=; пункты перечислений должны либо начинаться с заглавной буквы и заканчиваться точкой, либо начинаться со строчной и закачиваться точкой с запятой; последний пункт пере\-числений всегда заканчивается точкой.
%     \item Перед выкладыванием финальной версии необходимо просмотреть лог \LaTeX'a, и обратить внимание на сообщения вида \emph{Overfull \textbackslash hbox (1.29312pt too wide) in paragraph}. Обычно, это означает, что текст выползает за поля, и надо подсказать, как правильно разделять слова на слоги, чтобы перенос произошел автоматически.
%           Это делается, например, так: \verb=соломо\-волокуша=.
%     \item \emph{Обязательно используйте инструменты автоматической проверки правописания!}
%           Не посылайте текст даже научному руководителю на проверку, если не прогнали спеллчекер%
%           \footnote{Например, для браузера можно использовать плагин LanguageTool, для VS Code --- Code Spell Checker с расширением для русского (но он вроде не умеет пунктуацию, так что следите за запятыми).}%
%           , желательно, умеющий проверять пунктуацию~--- у научника куча времени уйдёт на комментарии вида \enquote{тут нужна запятая} и не останется сил посоветовать что-нибудь по делу.
%           А ещё научник будет шокирован уровнем грамотности современной молодёжи, впадёт в депрессию и не будет отвечать вам неделю.
% \end{itemize}

% \subsection{Листинги, картинка и прочий \enquote{не текст}}

% Различный \enquote{не текст} имеет свойство отображаться не там, где он написан в текстовом виде в \LaTeX{}, поэтому у него должна быть самодостаточ\-ная понятная подпись \verb=\caption{...}=, уникальная метка \verb=\label{...}=, чтобы на неё можно было бы ссылаться в тексте с помощью \verb=\ref{...}= (более того, ссылка из текста обязательна).
% Ниже Вы сможете увидеть таблицу \ref{time_cmp_obj_func}, на которую мы сослались буквально только что.

% \enquote{Не текста} может быть довольно много~--- чтобы не засорять корневую папку, хорошим решением будет складывать весь \enquote{не текст} в папку \texttt{figures}.
% Заклинание \verb=\includegraphics{}= уже знает этот путь и будет искать там файлы без указания папки.
% Команда \verb=\input{}= умеет ходит по путям, например \verb=\input{figures/my_awesome_table.tex}=.
% Кроме того, листинги кода можно подтягивать из файла с помощью команды \verb=\inputminted{file}=.

% %% Вставка кода с помощью listings
% % \begin{lstlisting}[caption={Название для листинга кода. Достаточно длинное, чтобы люди, которые смотрят картинку сразу после названия статьи (т.~е. все люди), смогли разобраться и понять к чему в статье листинги, картинки и прочий \enquote{не текст}.}, language=Caml, frame=single]
% %   let x = 5 in x + 1
% % \end{lstlisting}

% %% Вставка кода с помощью minted
% \begin{listing}
%     \caption{Название для листинга кода. Достаточно длинное, чтобы люди, которые смотрят картинку сразу после названия статьи (т.~е. все люди), смогли разобраться и понять к чему в статье листинги, картинки и прочий \enquote{не текст}.}
%     \begin{minted}[frame=single]{ocaml}
%     let x = 5 in x + 1
%   \end{minted}
% \end{listing}

% \subsubsection{Выделение куска листинга с помощью tikz}
% Это бывает полезно в текстах, а ещё чаще~--- в презентациях.
% Пример сделан на основе вопроса на \textsc{StackExchange}%
% \footnote{Вопрос про рамку вокруг листинга на StackExchange, \url{https://tex.stackexchange.com/questions/284311} (дата обращения: \DTMdate{2022-12-16}).}.
% Заодно тут показывается альтернативный minted пакет, lstlistings, если не хотите ставить Python и пакет pygments.
% В этом случае закомментируйте всё, что связано с minted, в matmex-diploma-custom.cls.
% И внимательно следите за тем, чтобы везде использовался либо только lstlistings, либо minted~--- смешивание их в одном документе приведёт к странным ошибкам.

% \begin{figure}
%     % TODO(Kakadu): Сделать \lstset глобально, чтобы не выписывать все опции листингов каждый раз
%     \begin{lstlisting}[ escapechar=!, keepspaces=true, extendedchars=\true, texcl=true
%                       , basicstyle=\ttfamily, commentstyle=\color{eclipseGreen}\ttfamily\itshape, language=c ]
% #include <stdio.h>
% #include <math.h>

% /** A comment in English */
% int main(void)
% {
%   double c = -1;
%   double z = 0;

%   // Это комментарий на русском языке
%   printf ("For c = %lf:\n", c);
%   for (int i = 0; i < 10; i++ ) {
%     printf( !\tikzmark{a}!"z %d = %lf\n"!\tikzmark{b}!, i, z);
%     z = pow(z, 2) + c;
%   }
% }
% \end{lstlisting}

%     \begin{tikzpicture}[use tikzmark]
%         \draw[fill=gray,opacity=0.1]
%         ([shift={(-3pt,2ex)}]pic cs:a)
%         rectangle
%         ([shift={(3pt,-0.65ex)}]pic cs:b);
%     \end{tikzpicture}
%     \caption{Пример листинга c помощью пакета \texttt{listings} и \textsc{TIKZ} декорации к нему, оформленные в окружении \texttt{figure}.
%         Обратите внимание, что рисунок отображается не там, где он в документе, а может \enquote{плавать}.}
% \end{figure}

% \subsection{Некоторые детали описания реализации}
% Описание реализации~--- очень важный раздел для будущих программных инженеров, т.е. почти для всех.
% Важно иметь всегда, даже если Вы написали прототип на коленке или немного скриптов.

% В процессе работы можно сделать огромное количество косяков, неполный список которых ниже.

% \begin{enumerate}
%     \item Реализация должна быть.
%           На публично доступную реализацию обязательна ссылка (в заключении, но можно продублировать тут).
%           Если код под \textsc{NDA}, то об этом, во-первых, должно быть сказано явно,
%           во-вторых, на защиту должны выно\-ситься другие результаты (например, архитектура), чтобы комис\-сия имела возможность оценить хоть что-то,
%           и, в третьих, должна быть справка от работодателя, что Вы правда что-то сделали.
%           \begin{itemize}
%               \item Рецензент обязан оценить код (о возможности должен побеспо\-коиться обучающийся).
%           \end{itemize}
%     \item Код реализации должен быть написан защищающимся целиком.
%           \begin{itemize}
%               \item Если проект групповой, то нужно явно выделить, какие части были модифицированы защищающимся.
%                     Например, в преды\-дущих разделах на картинке архитектуры нужно выделить цветом то, что Вы модифицировали.
%               \item Нельзя пускать в негрупповой проект коммиты от других людей, или людей не похожих на Вас.
%                     Например, в 2022 году защищающийся-парень делал коммиты от сценического псев\-донима, который намекает на женский \enquote{гендер}.
%                     (Нет, это не шутка.)
%                     На тот момент в российской культуре это выглядело странно.
%               \item Возможна ситуация, что вы используете конкретный ник в интернете уже лет пять, и желаете писать ВКР под этим ником на \GitHub{}.
%                     В принципе, это допустимо, но если Вы встретите преподавателя, который считает наоборот, то Вам придется грамотно отмазы\-ваться.
%                     В Вашу пользу могут сыграть те факты, что к нику на \GitHub{} у Вас приписаны настоящие имя и фамилия; что в репозитории у вас видна домашка за первый курс;
%                     и что Ваш преподаватель практики сможет подтвердить, что Вы уже несколько лет используете это ник; и т.п.
%           \end{itemize}
%     \item Если Вы получаете диплом о присвоении квалификации программиста, код должен соответствовать.
%           \begin{enumerate}
%               \item Не стоит выкладывать код одним коммитом.
%               \item Не стоит выкладывать код аккурат перед защитой.
%               \item Лучше хоть какие-то тесты, чем совсем без них.
%                     В идеале нужно предъявлять процент покрытия кода тестами.
%               \item Лучше сделать \textsc{CI}, а также \textsc{CD}, если оно уместно в Вашем проекте.
%               \item Не стоит демонстрировать на защите, что Вам даже не пришло в голову напустить на код линтеры и т.п.
%           \end{enumerate}
%     \item Если ваша реализация по сути является прохождением стандартного туториала,
%           например, по отделению картинок кружек от котиков с помощью машинного обучения, то необходимо срочно сообщить об этом руководителю практики/ВКР,
%           иначе Государственная Экзаменацион\-ная Комиссия \enquote{порвёт Вас как Тузик грелку}, поставит \enquote{единицу},
%           а все остальные Ваши сокурсники получат оценку выше.
%           (Это не шутка, а реальная история 2020 года.)
% \end{enumerate}

% \noindent Если Вам предстоит защищать учебную практику, а эти рекомендации видятся как более подходящие для защиты ВКР, то ... отмаза не засчиты\-вается, сразу учитесь делать нормально.
